\section{Inteligencia artificial y modelos de lenguaje}
\label{sec:ia}

La inteligencia artificial generativa constituye una parte importante del motor del proyecto. En esta sección se introducen los conceptos fundamentales sobre los modelos de lenguaje, el prompting engineering y los proveedores de API utilizados para acceder a los diferentes modelos.


\subsection{Modelos de lenguaje y IA generativa}
\label{subsec:llms}

Los modelos de lenguaje de gran tamaño (LLM, \emph{Large Language Models}) son sistemas de inteligencia artificial entrenados sobre grandes volúmenes de datos. Su arquitectura se basa en redes neuronales de tipo \emph{Transformer}~\cite{attention_is_all_you_need}, introducidas en 2017, que sustituyeron el procesamiento secuencial del texto por un mecanismo de atención capaz de analizar todas las palabras simultáneamente y capturar relaciones de largo alcance entre ellas.

A diferencia de los sistemas tradicionales de procesamiento de lenguaje natural, los LLM no se programan con reglas explícitas, sino que aprenden a predecir la siguiente unidad de texto, denominada token, a partir de las anteriores. Así es como se da lugar a modelos capaces de mantener conversaciones, redactar textos, traducir idiomas o como en este caso, generar código funcional a partir de instrucciones en lenguaje natural.
Cada modelo tiene una ventana de contexto llamada \textbf{context window}, que limita el número máximo de tokens que puede procesar simultáneamente, incluyendo tanto el prompt de entrada como la respuesta generada. Superar este límite implica que el modelo pierde información de los mensajes más antiguos. Este límite es muy importante, ya que los prompts de generación incluyen el esquema de datos, ejemplos del dataset e instrucciones de formato, lo que puede consumir varios miles de tokens por llamada, como se detalla en el capítulo~\ref{chap:diseno}.

Un fenómeno frecuente en los LLM son las alucinaciones, estas ocurren cuando el modelo genera contenido aparentemente valido pero en realidad es incorrecto o directamente inexistente~\cite{hallucination_survey}. A diferencia de un error de compilación, una alucinación no produce un fallo sino una salida que parece válida pero que contiene datos inventados, nombres de campos que no existen en el dataset o código erróneo. Las alucinaciones son el principal riesgo del sistema, ya que un campo inventado en el modelo Django generado provoca errores silenciosos en tiempo de ejecución. El diseño de los mecanismos de protección descritos en la Sección~\ref{subsec:proteccion-anti-alucionaciones} responden directamente a este problema.


La irrupción mediática de modelos como ChatGPT en 2022~\cite{chatgpt_launch} marcó 
el inicio de una etapa de adopción masiva de la IA generativa. Desde entonces han 
aparecido numerosos modelos, tanto comerciales (GPT-4, Claude, Gemini) como de 
pesos abiertos (LLaMA, Mistral, DeepSeek), accesibles a través de APIs públicas 
o ejecutables localmente. Esta diversidad permite a WebBuilder elegir el modelo 
más adecuado para cada tarea en función de criterios como coste, latencia o 
calidad de respuesta.

En este trabajo, los LLM no se utilizan como simples chatbots, sino como componentes programables dentro de un flujo automatizado: el modelo recibe datos estructurados, devuelve análisis razonados y produce código fuente que se integra directamente en el proyecto generado.


\subsection{Ingeniería de prompts}
\label{subsec:prompts}

La ingeniería de prompts (\emph{prompt engineering}) es la disciplina que estudia cómo formular las instrucciones enviadas a un modelo de lenguaje para obtener respuestas precisas y acorde a una necesidad, intentando a su vez reducir el numero de alucinaciones. A diferencia de la programación tradicional, donde el comportamiento del sistema se especifica mediante código, en los LLM el comportamiento se controla a través del lenguaje natural del propio prompt.

Para poder diseñar un buen prompt debemos seguir y combinar varios elementos: 

\begin{itemize}
    \item \textbf{Rol del modelo:} una descripción clara de la identidad o 
    función que debe asumir el modelo durante la tarea.
    \item \textbf{Instrucciones detalladas:} una especificación precisa de 
    lo que se espera que el modelo haga y cómo debe hacerlo.
    \item \textbf{Formato de respuesta:} indicaciones explícitas sobre la 
    estructura que debe tener la salida, como JSON, código Python o texto 
    plano.
    \item \textbf{Restricciones:} limitaciones explícitas que acotan el 
    comportamiento del modelo y reducen la probabilidad de respuestas 
    incorrectas o fuera de contexto.
    \item \textbf{Ejemplos de referencia:} muestras del estilo de salida 
    esperado que guían al modelo hacia el resultado deseado. Esta técnica, 
    conocida como \emph{few-shot prompting}, ha demostrado mejorar 
    notablemente la calidad de las respuestas en tareas complejas o con 
    formatos estrictos~\cite{openai_prompt_engineering}.
\end{itemize}


La calidad del resultado generado por un LLM depende, en \textbf{gran medida}, de la precisión del prompt. Un prompt ambiguo puede provocar respuestas inconsistentes, alucinaciones o salidas que no respetan el formato esperado. Por ello, el diseño, la iteración y el refinamiento de los prompts constituye una parte central del trabajo de desarrollo en sistemas basados en LLM, comparable en importancia a la propia escritura del código que los invoca.

La calidad del resultado no depende únicamente de la precisión del prompt, sino también de la capacidad del modelo utilizado. Un modelo con poco entrenamiento 
o potencia insuficiente producirá respuestas pobres independientemente de la calidad del prompt. Ambos factores, diseño del prompt y elección del modelo, son por tanto determinantes en el resultado final.

En el contexto de WebBuilder, la ingeniería de prompts juega un papel crítico en dos fases: el análisis de los datos de la API y la generación del código Django del proyecto. Cada uno de estos prompts ha sido diseñado tras innumerables pruebas para producir resultados coherentes y de calidad profesional, como se detallará en el Capítulo~\ref{chap:diseno}. Además, durante el proyecto se fueron guardando los datos de cada generación para realizar un estudio exhaustivo de qué funciona mejor para cada caso, cuyos resultados se recogen en el Capítulo~\ref{chap:experimentos}.


\subsection{Proveedores de modelos}
\label{subsec:proveedores}

Para acceder a un LLM y poder usarlo, en la mayoría de casos es necesario recurrir a un proveedor de API que disponga del modelo y ofrezca su uso. Cada vez existen mas proveedores en el mercado, cada uno con su propio catalogo. La gran mayoría de estos proveedores siguen el estándar de OpenAI descrito a continuación, lo que permite conectar de forma sencilla con WebBuilder. Cabe aclarar que los modelos gratuitos de los proveedores mencionados a continuación realizan bien su trabajo, pero nunca de forma perfecta. Las alucinaciones a día de hoy siguen existiendo sobre todo en los modelos gratuitos.

\textbf{Groq}~\cite{groq_docs} es un proveedor especializado en inferencia de baja latencia. Utiliza hardware propio (LPU, \emph{Language Processing Units}) optimizado, lo que le permite ofrecer velocidades de generación muy superiores a las de proveedores tradicionales basados en GPU. Groq aloja modelos de pesos abiertos como LLaMA, y dispone de un nivel de uso gratuito con límites de peticiones por minuto y por día. Su baja latencia hace que sea la opción preferente, donde la generación de templates es rapidísima.

\textbf{OpenRouter}~\cite{openrouter_docs} es un proveedor que da acceso a decenas de modelos de distintos proveedores a través de una única API. Lo mejor de este proveedor es la sencillez para la conexión: una sola clave de API permite consultar modelos de OpenAI, Anthropic, Google, Meta o cualquier otro proveedor integrado en la plataforma. Esto facilita las pruebas comparativas entre modelos y reduce el coste de cambiar de proveedor en caso necesario.

Otros proveedores como \textbf{NVIDIA NIM}~\cite{nvidia_nim_docs}, también compatibles con este estándar, amplían el catálogo con modelos especializados en código y razonamiento.

Los tres proveedores descritos comparten el denominado formato \textit{OpenAI-compatible}, un estándar surgido a partir de la API pública de OpenAI~\cite{openai_api_docs}. Este estándar define un único endpoint \texttt{/chat/completions} que acepta una lista de mensajes con roles diferenciados (\texttt{system}, \texttt{user} y \texttt{assistant}) y devuelve la respuesta del modelo en un formato JSON uniforme. Su adopción masiva por parte de los principales proveedores del mercado lo ha convertido en el protocolo común de la industria de modelos de lenguaje.

Esto es algo muy importante para el proyecto: al construir el cliente LLM sobre este protocolo común, el sistema puede cambiar de proveedor o de modelo modificando únicamente tres parámetros de configuración (la URL base, la clave de API y el identificador del modelo), sin necesidad de modificar ninguna línea de la lógica interna. Esta decisión de diseño, detallada en la Sección~\ref{subsec:conexion-llm}, es la que hace posible que WebBuilder soporte múltiples proveedores de forma nativa y que cualquier usuario pueda conectar su propio modelo personalizado.